\documentclass[conference]{IEEEtran}
\usepackage{cite}
\usepackage{amsmath,amssymb,amsfonts}
\usepackage{booktabs}
\usepackage{array}
\usepackage{url}
\usepackage{xcolor}
\usepackage{tabularx}
\usepackage{multirow}
\newcommand{\BREL}{\textsc{BREL}}
\newcommand{\MZM}{\textsc{MZM}}
\newcommand{\TQC}{\textsc{TQC}}
\newcommand{\AI}{\textsc{AI}}
\begin{document}

\title{Can Topological Quantum Computing Be Realized?\\
\large A Braiding-Readiness Evidence Ledger Research Proposal}

\author{\IEEEauthorblockN{Anonymous Research Proposal}
\IEEEauthorblockA{\textit{Evidence-to-Design Scientific Workflow}\\
Research Plan for a Falsifiable Topological-Quantum-Computing Primitive}}

\maketitle

\begin{abstract}
Topological quantum computing (TQC) offers an unusually compelling answer to the noise problem: encode quantum information nonlocally in non-Abelian anyons, manipulate it through braid operations, and make selected local perturbations unable to resolve the computational state. The theoretical mechanism is well developed, while solid-state Majorana-zero-mode (MZM) platforms have produced candidate phase and transport signatures. Those two facts do not yet imply that a topological quantum computer has been realized. A local zero-bias feature, a gap proxy, or a favorable isolated device trace cannot by itself establish nonlocal encoding, non-Abelian exchange, protected logical operation, or universal computation. This paper synthesizes the evidence and proposes the Braiding-Readiness Evidence Ledger (BREL), a pre-registered, conjunctive realization criterion with five columns: phase, nonlocality, operation, protection, and architecture. The proposed future study uses matched trivial controls, hold-out devices, blinded analysis, fusion/exchange order controls, scaling tests, and an explicit non-Clifford resource budget. Its primary endpoint is a staged platform classification rather than a selective Majorana label. The direct conclusion is that TQC is theoretically realizable and deserves serious physical implementation research; a platform should earn the stronger ``braiding-ready topological-qubit primitive'' claim only after it passes the linked BREL evidence chain. This is a design-only proposal and reports no fabricated devices, measurements, simulations, or observed qubit results.
\end{abstract}

\begin{IEEEkeywords}
topological quantum computing, non-Abelian anyons, Majorana zero modes, topological qubits, braiding, quantum error correction, condensed-matter experiment design, falsification.
\end{IEEEkeywords}

\section{Introduction}

Quantum information is valuable precisely because a coherent quantum state can represent and transform amplitudes that a classical bit cannot. It is also fragile. A useful quantum computer must keep its computational degrees of freedom sufficiently isolated from uncontrolled environmental interactions while still allowing deliberate initialization, gates, measurement, calibration, and scaling. This tension has shaped the major hardware programs in superconducting circuits, trapped ions, neutral atoms, photonics, spins, and topological matter. In each program, the question is not merely whether a quantum effect can be observed. It is whether that effect can be turned into an operational, reproducible, and fault-tolerant computing resource.

Topological quantum computing proposes a distinctive route. Its ideal information carrier is not a local microscopic two-level system whose state can be read at a point. Instead, information is encoded in a collective, spatially distributed fusion space of non-Abelian anyons. Exchanging identical quasiparticles along different braids acts on this space by matrices that need not commute. Under the assumptions of the topological phase, a local perturbation cannot readily distinguish the global encoded state. Kitaev formulated anyonic fault-tolerant computation in this language \cite{kitaev2003}; Nayak \textit{et al.}\ provide a canonical account of how non-Abelian degeneracy, braiding, and measurement form a computational model \cite{nayak2008}.

This prospect motivates the question in the title. It is tempting to answer it with a headline: a Majorana-like signature has appeared in a nanowire, therefore topological quantum computing has arrived; or a difficult experiment has not resolved every ambiguity, therefore the entire approach is impossible. Both inferences are poor science. The first confuses a candidate local signature with a computational primitive. The second ignores a principled theory and substantial materials/device progress. The right question is narrower and more useful: \emph{what linked experimental and engineering evidence is sufficient to establish that a specified platform has realized a controllable, protected topological-qubit primitive rather than a local mimic?}

This proposal answers with the \emph{Braiding-Readiness Evidence Ledger} (\BREL). BREL is a research architecture, not a new particle and not a claim of completed hardware. It connects five propositions that are frequently collapsed into one: a gapped candidate phase exists; the low-energy degree of freedom is nonlocal; a controlled fusion or exchange operation has a discriminating order dependence; the encoding suppresses a relevant error channel with measurable scaling; and the device can enter an explicit universal-computing architecture. A platform earns a strong realization label only when all five propositions pass predeclared controls. A partial result remains informative because it identifies the first unsupported link.

The paper makes four contributions. First, it corrects common conceptual shortcuts about topology, Majorana modes, temperature, and error correction. Second, it turns the broad feasibility question into a hierarchy of falsifiable claims. Third, it proposes a D0--D4 future study design that joins condensed-matter measurements to quantum-information requirements. Fourth, it defines a formal decision rule and an evidence-governance model suitable for a multi-device research program. The paper is a proposal: it authorizes no cryogenic operation, fabrication, data acquisition, hardware control, or numerical experiment, and it reports no observed result.

\section{Survey: What Is Established and What Is Not}

\subsection{Topology is an information-protection mechanism, not a geometric slogan}

The topological approach is often explained through a picture of particle paths. The picture is useful only if its physical content remains precise. In two spatial dimensions, paths of identical particles can form braids rather than merely exchange labels. For non-Abelian anyons, the final quantum transformation depends on the braid history and acts on a degenerate multi-particle state space. The relevant states are emergent quasiparticle states in a many-body system; they are not ordinary particles changing between a conventional manifold and a Mobius strip. Their non-Abelian exchange statistics are the computational resource \cite{nayak2008}.

In a Majorana realization, idealized zero modes obey the Majorana relation
\begin{equation}
\gamma_j=\gamma_j^\dagger,\qquad \{\gamma_i,\gamma_j\}=2\delta_{ij},
\label{eq:majorana}
\end{equation}
where $\gamma_j$ is a Majorana operator and $\delta_{ij}$ is the Kronecker delta. A pair can define a fermionic mode, while a suitably constrained collection of separated modes supplies a parity/fusion space. An exchange of modes $i$ and $j$ is ideally represented, up to a convention-dependent phase, by
\begin{equation}
U_{ij}=\exp\left(\frac{\pi}{4}\gamma_i\gamma_j\right).
\label{eq:braid}
\end{equation}
The physical implication of \eqref{eq:braid} is not that an arbitrary endpoint peak is a braid. It is that a multi-mode device must support an order-sensitive operation and measurement with a well-defined relation to the encoded state.

The proposed protection is equally specific. Nonlocal encoding can suppress the sensitivity of selected logical degrees of freedom to local perturbations. It does not eliminate all error sources. Finite temperature, quasiparticle poisoning, disorder, finite-size hybridization, nonadiabatic control, readout error, classical control error, and architectural overhead remain. A serious TQC program therefore does not promise to dispense with cooling or error correction. It seeks a physical encoding that changes the error model and can reduce the burden imposed on higher-level fault tolerance.

\subsection{Candidate platforms and the Majorana evidence boundary}

Semiconductor-superconductor heterostructures are a major solid-state candidate. Effective one-dimensional models predict a topological regime when spin-orbit coupling, Zeeman splitting, chemical potential, and proximity-induced pairing jointly meet stated conditions \cite{lutchyn2010,oreg2010}. A common idealized phase-boundary expression is
\begin{equation}
V_Z^2>\mu^2+\Delta^2,
\label{eq:phaseboundary}
\end{equation}
where $V_Z$ is an effective Zeeman energy, $\mu$ is chemical potential, and $\Delta$ is induced pairing. Equation \eqref{eq:phaseboundary} is a model-level guide, not a device verdict: disorder, multiple bands, inhomogeneous electrostatics, finite length, and imperfect interfaces can change the observable phenomenology.

Lutchyn \textit{et al.}\ review the materials progress in InAs/InSb semiconductor-superconductor structures, tunneling and Coulomb-blockade signatures, and the importance of future fusion and exchange experiments \cite{lutchyn2018}. Albrecht \textit{et al.}\ reported exponential protection behavior in Majorana islands \cite{albrecht2016}; Zhang \textit{et al.}\ reported quantized Majorana conductance in hybrid devices \cite{zhang2018}; and Vaitiekenas \textit{et al.}\ investigated flux-induced topological superconductivity in full-shell nanowires \cite{vaitiekenas2020}. These advances make a physical realization program credible. They do not establish non-Abelian operations or a general-purpose topological computer.

The difference is consequential because local low-energy signals can have non-topological explanations. Pan and Das Sarma show that trivial Andreev-bound-state mechanisms can produce zero-bias conductance peaks in Majorana-nanowire settings \cite{pan2020}. A local spectral anomaly may be necessary, suggestive, or useful for screening; it is not a unique identifier of a non-Abelian degree of freedom. This proposal treats that ambiguity as a design requirement rather than an excuse for paralysis: every strong candidate claim must face a matched trivial control and an independent nonlocal/operational test.

\subsection{From protected operations to useful computation}

Theory also imposes an architectural boundary. Majorana/Ising braiding natively supplies a protected Clifford-type operation set, but Clifford gates alone are not universal for quantum computation. A complete program must identify the non-Clifford resource, its preparation/measurement path, its error contribution, and how it enters a logical resource budget \cite{bravyi2005}. Aasen \textit{et al.}\ identify milestones toward Majorana-based computation, and Karzig \textit{et al.}\ describe scalable architectures designed to protect against quasiparticle poisoning \cite{aasen2016,karzig2017}. Those works motivate a staged approach: phase evidence, nonlocal encoding, measurement/operation primitives, scaling, and an architecture-level error budget.

Table~\ref{tab:levels} separates the evidence levels. The table is intentionally asymmetric. An upper row may be necessary for a lower row, but a lower row is not inferred automatically from any upper row. This avoids granting a universal computational label to a platform that has only shown a physical precursor.

\begin{table*}[t]
\caption{Evidence hierarchy for a topological-quantum-computing claim}
\label{tab:levels}
\centering
\footnotesize
\begin{tabular}{p{0.18\textwidth} p{0.31\textwidth} p{0.38\textwidth}}
\toprule
\textbf{Evidence level} & \textbf{What it can support} & \textbf{What it cannot support by itself}\\
\midrule
Candidate phase & A reproducible gapped parameter region compatible with a model of topological superconductivity. & Nonlocal encoded information, non-Abelian exchange, logical operations, or a topological computer.\\
Local signature & A low-energy feature consistent with an MZM hypothesis and a reason for further screening. & Exclusion of a trivial Andreev-bound-state explanation or proof of a braiding algebra.\\
Nonlocality/parity structure & A spatially distributed degree of freedom beyond a single-end signal, conditional on null controls. & A controlled computational operation or protected logical scaling.\\
Fusion/exchange operation & Evidence for a specified order-dependent operation when matched controls are rejected. & Multi-qubit scalability, universal computation, or a fault-tolerance threshold.\\
Protection and architecture & A candidate primitive with an explicit error/resource path to useful computation. & General performance claims outside the specified device, protocol, and operating regime.\\
\bottomrule
\end{tabular}
\end{table*}

\section{Research Gaps, Questions, and Hypotheses}

The survey yields five connected gaps. First, the field lacks a generally transparent realization rule that joins material evidence to the computational claim. Second, local MZM-like signatures can be non-diagnostic. Third, a favorable fusion or braid-like signal is not automatically an engineering demonstration of protected operation. Fourth, the protected-to-universal bridge is often left implicit even though a non-Clifford resource is essential. Fifth, cross-device reproducibility and falsification controls are too easily secondary to a favorable example trace.

The proposed study asks the following questions:
\begin{enumerate}
\item Can a pre-registered multi-axis ledger separate a candidate topological-qubit primitive from matched local-null mechanisms more reliably than a single-signature workflow?
\item Does a specified device network show nonlocal parity/fusion behavior and exchange-order dependence that persist across held-out devices and are absent in deliberately trivial controls?
\item Does the proposed encoding show a protection-relevant scaling relation with separated-encoding design variables after accounting for ordinary control and readout effects?
\item Can the resulting primitive enter an explicit protected-Clifford plus non-Clifford resource budget with a favorable, auditable condition for future logical computation?
\end{enumerate}

The central hypothesis is deliberately conditional:
\begin{quote}
\emph{If a gate-defined candidate network contains a controllable topological Majorana sector, then phase, nonlocality, fusion/exchange order dependence, protection scaling, and architecture accounting will form a coherent pattern that the preregistered local-null families and matched trivial controls cannot jointly reproduce.}
\end{quote}

This claim has direct falsifiers. The hypothesis is rejected or paused if a trivial/scrambled control reproduces the full pattern; if reversing the operational order creates no contrast beyond readout bias and drift; if remote detuning lacks the predicted nonlocal consequence; if apparent protection is explained by a local covariate; or if no credible non-Clifford resource path survives the error budget. These are not merely caveats. They are the scientific tests that make an affirmative result meaningful.

\section{The Braiding-Readiness Evidence Ledger}

\subsection{A conjunctive architecture for a conjunctive claim}

The selected idea is the \BREL. It is an experimental-decision architecture rather than a weighted score that could conceal a failed prerequisite beneath high values elsewhere. The five columns are:
\begin{enumerate}
\item \textbf{Phase:} a candidate gapped operating region is stable and reproducible under the declared conditions.
\item \textbf{Nonlocality:} a paired-end or parity structure produces remote-correlated behavior that the matched local null cannot explain.
\item \textbf{Operation:} a predeclared fusion/exchange protocol shows the expected order-dependent outcome relative to matched order and trivial controls.
\item \textbf{Protection:} an error or instability proxy displays the required scaling with a separated-encoding design variable after accounting for confounds.
\item \textbf{Architecture:} protected operations, measurement, poisoning/readout contributions, and the non-Clifford resource are expressed in a traceable universal-computation budget.
\end{enumerate}

Figure~\ref{fig:brel} shows the logic. The proposed workflow lets a future research team stop early when a prerequisite fails. This is more efficient and more honest than spending scarce measurement time trying to salvage a universal claim from an ambiguous local signal. It also makes a partial success actionable: a platform may be an interesting phase candidate but not yet a nonlocal candidate, or it may have an operation signal whose protection scaling is inadequate.

\begin{figure*}[t]
\centering
\fbox{\begin{minipage}{0.92\textwidth}
\centering
\vspace{3pt}
\textbf{Candidate device/network} $\rightarrow$ \textbf{D1 phase screen} $\rightarrow$ \textbf{D2 nonlocality/parity} $\rightarrow$ \textbf{D3 fusion/exchange order test} $\rightarrow$ \textbf{D4 protection and architecture}\\[4pt]
\small Each stage is evaluated against a preregistered candidate model, matched trivial control, quality-control record, and held-out device/batch condition.\\[5pt]
\textbf{Fail or inconclusive} $\rightarrow$ \textbf{diagnostic classification and redesign} \quad $|$ \quad \textbf{all five columns pass} $\rightarrow$ \textbf{braiding-ready primitive claim}\\[-1pt]
\small No stage alone establishes a general-purpose topological quantum computer. The present paper authorizes only the design, not any physical execution.
\vspace{3pt}
\end{minipage}}
\caption{Braiding-Readiness Evidence Ledger (BREL). The diagram is a proposed research decision structure, not a report of device data or a completed experiment.}
\label{fig:brel}
\end{figure*}

\subsection{Ledger evidence and null controls}

Table~\ref{tab:brel} specifies the independent evidence required in each column. The central methodological move is to make the alternative explanation operational. A local Andreev-bound-state family, a disorder/multiband subgap-state family, and a pulse/readout-artifact family are entered into the preregistration before any confirmatory conclusion. A control is not an unrelated poor device. It is a carefully matched configuration that shares as much geometry, pulse timing, readout, and analysis as possible while removing or scrambling the hypothesized topological condition.

The same principle applies to evidence selection. A visually striking trace is not an endpoint. The future study reports all eligible devices under a predeclared inclusion rule, device and batch provenance, calibration status, and the complete distribution of the target statistic. It separates exploratory diagnostics from confirmatory analysis. This protects against the common failure mode in which a protocol is changed until an expected feature appears and is then reported as though it were a preregistered test.

\begin{table*}[t]
\caption{BREL columns, observables, controls, and claim boundary}
\label{tab:brel}
\centering
\footnotesize
\begin{tabular}{p{0.12\textwidth} p{0.25\textwidth} p{0.25\textwidth} p{0.25\textwidth}}
\toprule
\textbf{Column} & \textbf{Future target evidence} & \textbf{Matched falsification control} & \textbf{Permitted conclusion if passed}\\
\midrule
Phase & Reproducible gap-proxy map and candidate region across the declared device/batch set. & Intentionally trivial configuration, predeclared nuisance ranges, and held-out batch. & Candidate phase supported for the named platform and regime.\\
Nonlocality & Paired-end/parity statistic and remote-control response outside the local-null envelope. & Local-state family fitted under the same protocol; remote segment detuned/scrambled. & Nonlocal candidate supported, conditional on the defined null family.\\
Operation & Reproducible contrast between protocol orders in fusion/parity outcomes with calibrated readout. & Order reversal/scramble, matched trivial network, and pulse-artifact control. & Specified order-dependent operation supported; not yet a scalable computer.\\
Protection & Error/instability proxy improves with separation, exposure, time, or temperature in the predicted direction. & Local-covariate model, matched readout control, and independent calibration reference. & Evidence that the specified encoding suppresses a defined error channel.\\
Architecture & Versioned protected-Clifford, measurement, poisoning, readout, and non-Clifford resource budget. & Sensitivity analysis against adverse but plausible component errors. & A conditional route to useful logical computation is specified.\\
\bottomrule
\end{tabular}
\end{table*}

\subsection{Realization classification}

BREL produces four allowable classifications. \emph{Not supported} means at least one prerequisite is inconsistent with the predeclared target after quality checks. \emph{Phase candidate} means the first column alone passes. \emph{Nonlocal candidate} requires phase and nonlocality. \emph{Braiding-ready topological-qubit primitive} requires every ledger column, including a universal-resource accounting condition, to pass. A more expansive claim---a scalable topological quantum computer---requires additional evidence beyond this proposal: multiple encoded qubits, logical operation demonstrations, a platform-specific fault-tolerance analysis, and system-level scaling.

This language is intentionally neither promotional nor defeatist. It permits a precise affirmative result when the evidence earns it. It prohibits a promotion of a single device feature into a claim that has not yet been tested.

\section{Proposed Study Design and Methods}

\subsection{Scope, unit of analysis, and execution boundary}

The future physical object is a gate-defined network of proximitized semiconductor-superconductor islands selected by a qualified device team as a candidate MZM platform. The unit of analysis is a device-protocol episode: a named device, batch, configuration version, environmental/calibration state, protocol order, readout record, and predeclared endpoint. The future study follows five stages, D0 through D4. The present proposal performs none of them. No hardware is fabricated, cooled, gated, pulsed, measured, simulated, or controlled; no data record has been accessed; and no outcome is reported.

Any physical execution is conditional on laboratory safety approval, device/facility procedures, instrument review, institutional research governance, data-management rules, and an expert assessment of whether the selected platform and protocol are appropriate. The design intentionally avoids fabrication recipes, operating set points, or automation instructions. Its contribution is the evidence architecture and analytical contract.

\subsection{D0: pre-registration and synthetic-only analysis validation}

D0 freezes the decision framework before future device data are interpreted. The future team specifies an effective candidate model, its stated parameter domain, the local-null model families, the candidate and trivial-control definitions, path orders, eligibility/exclusion rules, endpoint estimators, uncertainty intervals, multiplicity treatment, calibration records, and the unblinding plan. A future digital twin may test whether an analysis pipeline can distinguish the declared model families under hypothetical nuisance conditions. It must remain synthetic-only at this stage and must not be tuned on confirmatory device data.

D0 also establishes a blinded label key. Candidate/control labels, protocol order labels, and the planned hold-out batch are isolated from the primary analysis until the code and thresholds are frozen. The purpose is not bureaucratic ceremony. It prevents the same data from choosing the signal, choosing the threshold, and declaring the result. The D0 gate blocks any physical-data conclusion until the ledger and its alternative models are versioned.

\subsection{D1: phase and reproducibility screening}

D1 asks whether a candidate gapped regime is reproducible for the named platform. Independent variables include device class, candidate versus intentionally trivial configuration, declared control trajectory, device batch, and facility/session metadata. The future dependent variables include a gap-proxy map, the stable region under the declared criterion, candidate/control separation, and device-to-device yield. A common quality-control rule applies to all configurations.

The analysis reports the full eligible device set. When adequate data exist, a hierarchical model represents batch, device, and session variation rather than treating every trace as independent. The D1 gate passes only if the candidate criterion is met in the held-out batch and if a matched trivial configuration does not reproduce it within the preregistered practical-equivalence bound. Passing D1 supports a candidate phase in a stated operating regime. It does not support non-Abelian statistics.

\subsection{D2: nonlocality and parity structure}

D2 tests whether the relevant degree of freedom is spatially distributed in a way that a local model cannot explain. The future observables include paired-end correlation, remote-control response, parity-readout consistency, and a controlled separation/length dependence. The candidate and control configurations receive the same calibration, acquisition policy, and analysis code. The local-null envelope is fitted before the final confirmatory comparison.

The central test is relational: detuning or modifying the remote segment should have a predeclared consequence for the joint endpoint if the encoded degree of freedom is nonlocal. A positive outcome requires that result to recur in the held-out device/batch condition and to lie outside the local-null envelope with the declared uncertainty treatment. If D2 fails, the result is not discarded. The platform is classified as phase candidate at most, and the experimental team can decide whether the limitation comes from material quality, network design, readout, or an incorrect physical hypothesis.

\subsection{D3: fusion/exchange order-dependence}

D3 is the operational center of the study. It compares a future proposed fusion or exchange pathway with an order-reversed or scrambled pathway that is matched in total pulses, timing, and readout as far as the hardware permits. The primary future endpoint is a predeclared contrast in a fusion/parity outcome distribution. Secondary endpoints include repeatability, readout-bias calibration, device/batch variation, and the contrast observed in an intentionally trivial network.

Non-Abelianity is about the order of operations. Consequently, repeated use of one favorable pulse sequence is not a substitute for the ordered comparison. The protocol record must store configuration version, timing order, calibration status, readout method, and any intervention that occurred before unblinding. The D3 gate passes only if the target order contrast is reproduced in the stated held-out devices and is not recreated by the matched trivial and pulse-artifact controls. A positive D3 result would support the specified operation; it would still need D4 before it could be called a braiding-ready primitive.

\subsection{D4: protection scaling and universal-resource accounting}

D4 connects physics to computing value. The future study evaluates an error or instability proxy as a function of separation/length class, temperature condition, operation or exposure duration, and readout protocol. It models ordinary covariates such as device disorder proxy, finite-size hybridization, electrostatic drift, and readout fidelity. The target is a predeclared relation consistent with the purported protection mechanism, not simply a small error at one convenient point.

The architecture package then records which operations are protected, which are merely controlled, the sources of physical error, conversion assumptions from physical to logical error, and the proposed non-Clifford resource. The latter may be a magic-state or equivalent supplemental resource, but its preparation and error contribution cannot be omitted. The D4 gate passes only if protection scaling survives sensitivity to the null models and if the complete resource budget meets a future platform-specific feasibility condition. This plan does not invent that condition; a qualified architecture team must select and justify it before execution.

\section{Analysis, Decision Logic, and Expected Outcome Branches}

\subsection{Conjunctive primary endpoint}

Let $P$, $N$, $B$, $S$, and $U$ be preregistered binary indicators for the phase, nonlocality, braiding/fusion operation, protection scaling, and universal-resource columns, respectively. The confirmatory realization endpoint is
\begin{equation}
R_{\mathrm{BREL}}=P\,N\,B\,S\,U.
\label{eq:brel}
\end{equation}
Equation \eqref{eq:brel} forbids compensating a failed nonlocality test with an attractive phase map or a favorable architecture calculation. For each column $j$, the future protocol defines an estimated target effect $\widehat{T}_j$, a matched-control effect $\widehat{T}^{\mathrm{ctrl}}_j$, an admissible direction/range $\mathcal{A}_j$, a practical separation $\delta_j$, a quality-control indicator $Q_j$, and a hold-out reproducibility indicator $H_j$. The generic gate is
\begin{equation}
G_j=\mathbb{I}\left\{\widehat{T}_j\in\mathcal{A}_j,\;\left|\widehat{T}_j-\widehat{T}^{\mathrm{ctrl}}_j\right|>\delta_j,\;Q_j=1,\;H_j=1\right\}.
\label{eq:gate}
\end{equation}
Here $\mathbb{I}\{\cdot\}$ is the indicator function. Values for $\delta_j$, the estimators, and the interval procedure depend on the final physical platform; they are intentionally left for future expert pre-registration rather than invented after the fact.

The analysis gives effect sizes, compatible uncertainty intervals, full eligible-device distributions, and the result of each null-control comparison. If the data structure supports it, hierarchical models should represent batch, device, and session variation. The output must distinguish confirmatory tests from exploratory diagnostics. The study should not replace all evidence with a composite number: \eqref{eq:brel} is a classification rule, while the five column-wise records preserve scientific detail.

\subsection{Outcome branches}

Table~\ref{tab:branches} makes expected outcomes conditional. This is not a timid substitute for a conclusion; it states exactly what each result would support. A positive result in every column is a strong, publishable case for a narrowly stated primitive. A negative result has equal design value because it limits the physical hypothesis or identifies an engineering bottleneck.

\begin{table*}[t]
\caption{Future outcome branches and permitted conclusions}
\label{tab:branches}
\centering
\footnotesize
\begin{tabular}{p{0.25\textwidth} p{0.34\textwidth} p{0.28\textwidth}}
\toprule
\textbf{Observed future ledger pattern} & \textbf{Permitted interpretation} & \textbf{Next action}\\
\midrule
Phase fails or trivial control matches & The selected platform/configuration does not support the claimed candidate phase under the stated criterion. & Diagnose material/device model; do not proceed to a TQC claim.\\
Phase passes, nonlocality fails & A candidate phase/local signature is insufficient for a nonlocal encoded-state claim. & Revise network/readout or test alternative null mechanisms.\\
Phase and nonlocality pass, operation fails & A nonlocal candidate may exist, but controlled non-Abelian operation is not supported. & Audit coherence, path implementation, readout, and control design.\\
Operation passes, protection or architecture fails & The operation is promising, but computing advantage/fault-tolerant relevance is not established. & Improve scaling, poisoning/readout budget, or universal-resource plan.\\
All columns pass against controls and hold-outs & The named platform supports a braiding-ready topological-qubit primitive under its declared scope. & Replicate independently; expand to encoded multi-qubit/logical-operation studies.\\
\bottomrule
\end{tabular}
\end{table*}

\subsection{Reproducibility and evidence governance}

The project archives future raw data under the host laboratory's policy together with instrument/configuration metadata, calibration records, device batch identifiers, analysis environment versions, preregistration versions, and the unblinding record. The device inclusion/exclusion rule and all quality-control codes remain attached to the final table. At least one batch is held out until the criterion is fixed. When institutional and intellectual-property rules permit, the future team should release a non-sensitive analysis specification, summary-level evidence, and code sufficient to reproduce the ledger classification.

This governance is not a claim that the entire literature is wrong or that data sharing alone solves a complex physical inference. It is a safeguard against avoidable self-confirmation. The ledger requires the same discipline for candidate and control devices. A result that needs a different exclusion rule or analysis method only for the control is evidence of a workflow problem, not a clean pass.

\section{Risks, Human Review, and Evidence Limits}

The primary scientific risk is a false-positive topological claim caused by a local mimic. BREL addresses it through matched trivial configurations, multiple null model families, nonlocal endpoints, operation-order controls, and a conjunctive classification. The opposing risk is a false negative caused by insufficient readout fidelity, calibration loss, or environmental drift. The design therefore permits an explicit \emph{inconclusive} diagnostic state in development even though it does not permit an inconclusive result to be upgraded to a confirmatory pass.

Several reviews are indispensable before any physical execution. A condensed-matter theory review must assess whether the candidate and null families are adequate for the chosen device. An experimental review must assess the realism of the control network, measurement chain, and calibration plan. A statistical review must assess the primary estimators, practical separations, model sensitivity, and hold-out structure. A quantum-information architecture review must test the physical-to-logical error mapping and the non-Clifford resource budget. Finally, laboratory and facility review must govern all cryogenic, electrical, magnetic, and equipment procedures; this design itself performs none of those operations.

The source base also has limits. The Survey Agent attempted three focused dual-engine literature searches through OpenAlex and AnySearch, but both sources were unreachable from the local environment and returned no candidates because the network actively refused the connection. The report therefore does not claim a live, exhaustive literature census. Instead, it relies on registered cornerstone sources, with publisher metadata and abstracts checked directly for the Nayak \textit{et al.}\ review and the Lutchyn \textit{et al.}\ platform review. Full-text applicability, current device claims, and any claimed breakthrough require human source review before a real program begins. This limitation does not change the core theoretical distinction or the design requirement for controls; it prevents an overstatement of search completeness.

\section{Conclusion}

Topological quantum computing can be realized in the theoretical sense: non-Abelian anyon models provide a well-defined, fault-tolerant computational mechanism, and Majorana-based solid-state platforms give it a plausible physical route. The present state of evidence should not be compressed into a claim that a general topological computer already exists, nor into the opposite claim that difficult engineering makes the concept empty. The key task is a disciplined transition from suggestive platform evidence to an operational, protected primitive.

The proposed BREL study makes that transition falsifiable. It requires a reproducible phase, nonlocal encoding evidence, order-dependent fusion/exchange behavior, protection scaling, and a transparent universality budget, each tested against predeclared null controls and held-out devices. If the chain passes, the field gains a strong and precisely scoped realization claim. If it fails, the field gains a localized explanation for where the physical or engineering bridge broke. Either result advances the question more effectively than a single-signature search.

\appendices
\section{Idea Evolution and Direction Selection Audit}

The Idea Agent began with a broad seed: build a more sensitive Majorana detector. It rejected that seed as a primary contribution because a more detailed local signature does not resolve the local-mimic problem. A second candidate, a twin-device nonlocal-correlator benchmark, substantially improved the problem framing but did not force an operational exchange test or an architecture-level universal-resource condition. A third candidate, a protected-Clifford plus magic-state budget, addressed universality but could be disconnected from a physical phase claim. A high-risk interferometric non-Abelian witness was retained as a conditional future route because its coherence and calibration demands can make a null ambiguous.

BREL was selected because it combines the discriminating strength of nonlocal controls, the operational demand of path-order comparison, and the architecture discipline of a non-Clifford budget. Its rejection conditions are traceable to the Survey gap ledger: it makes a realization criterion explicit, treats local signatures as insufficient, prevents operational and engineering evidence from being conflated, exposes the universality bridge, and requires cross-device falsification.

\section{Evidence Coverage and Required Review}

The Author Agent used only sources registered by the Survey evidence plan. The references support theory, candidate platform status, local-mimic risk, architectural milestones, and the limitation of Majorana/Ising operations. The report does not cite a news announcement as experimental proof and does not declare a currently operating topological computer. The following items remain for qualified human review before any downstream research is authorized: full-text source applicability; device-specific material physics; laboratory safety and protocol feasibility; statistical thresholds; the validity range of all null model families; and the assumed non-Clifford-resource overhead.

\begin{thebibliography}{99}

\bibitem{kitaev2003}
A. Y. Kitaev, ``Fault-tolerant quantum computation by anyons,'' \emph{Annals of Physics}, vol. 303, no. 1, pp. 2--30, 2003, doi: 10.1016/S0003-4916(02)00018-0.

\bibitem{nayak2008}
C. Nayak, S. H. Simon, A. Stern, M. Freedman, and S. Das Sarma, ``Non-Abelian anyons and topological quantum computation,'' \emph{Reviews of Modern Physics}, vol. 80, pp. 1083--1159, 2008, doi: 10.1103/RevModPhys.80.1083.

\bibitem{lutchyn2010}
R. M. Lutchyn, J. D. Sau, and S. Das Sarma, ``Majorana fermions and a topological phase transition in semiconductor-superconductor heterostructures,'' \emph{Physical Review Letters}, vol. 105, 077001, 2010, doi: 10.1103/PhysRevLett.105.077001.

\bibitem{oreg2010}
Y. Oreg, G. Refael, and F. von Oppen, ``Helical liquids and Majorana bound states in quantum wires,'' \emph{Physical Review Letters}, vol. 105, 177002, 2010, doi: 10.1103/PhysRevLett.105.177002.

\bibitem{lutchyn2018}
R. M. Lutchyn, E. P. A. M. Bakkers, L. P. Kouwenhoven, P. Krogstrup, C. M. Marcus, and Y. Oreg, ``Majorana zero modes in superconductor-semiconductor heterostructures,'' \emph{Nature Reviews Materials}, vol. 3, pp. 52--68, 2018, doi: 10.1038/s41578-018-0003-1.

\bibitem{albrecht2016}
S. M. Albrecht \emph{et al.}, ``Exponential protection of zero modes in Majorana islands,'' \emph{Nature}, vol. 531, pp. 206--209, 2016, doi: 10.1038/nature17162.

\bibitem{zhang2018}
H. Zhang \emph{et al.}, ``Quantized Majorana conductance,'' \emph{Nature}, vol. 556, pp. 74--79, 2018, doi: 10.1038/nature26142.

\bibitem{vaitiekenas2020}
S. Vaitiekenas \emph{et al.}, ``Flux-induced topological superconductivity in full-shell nanowires,'' \emph{Science}, vol. 367, eaav3392, 2020, doi: 10.1126/science.aav3392.

\bibitem{pan2020}
H. Pan and S. Das Sarma, ``Physical mechanisms for zero-bias conductance peaks in Majorana nanowires,'' \emph{Physical Review Research}, vol. 2, 013377, 2020, doi: 10.1103/PhysRevResearch.2.013377.

\bibitem{aasen2016}
D. Aasen \emph{et al.}, ``Milestones toward Majorana-based quantum computing,'' \emph{Physical Review X}, vol. 6, 031016, 2016, doi: 10.1103/PhysRevX.6.031016.

\bibitem{karzig2017}
T. Karzig \emph{et al.}, ``Scalable designs for quasiparticle-poisoning-protected topological quantum computation with Majorana zero modes,'' \emph{Physical Review B}, vol. 95, 235305, 2017, doi: 10.1103/PhysRevB.95.235305.

\bibitem{bravyi2005}
S. Bravyi, ``Universal quantum computation with ideal Clifford gates and noisy ancillas,'' \emph{Physical Review A}, vol. 71, 022316, 2005, doi: 10.1103/PhysRevA.71.022316.

\end{thebibliography}

\end{document}
